\subsection{Neural Network Structure}
The dynamical systems investigated in this work were approximated using a feedforward neural network implemented in PyTorch. The objective of the neural network was to learn the relationship between time and the corresponding physical state of each system, without having explicit information about the differential equations governing these systems during training.

For systems with a single state variable, the network receives time $(t)$ as input and predicts the corresponding position:
\begin{equation}
x(t)=NN(t)
\end{equation}

where $x(t)$ represents the approximation of the physical solution obtained by the neural network.

For the double pendulum system, however, the output dimension is larger since the network predicts two angular positions:
\begin{equation}
(\theta_1(t),\theta_2(t))=NN(t)
\end{equation}

The general architecture, although modified for some specific experiments, consists of three fully connected layers with 128 neurons each. The activation function used between layers is the hyperbolic tangent function (Tanh), allowing the network to represent the nonlinear relationships present in the physical systems studied.

Additionally, the network parameters were optimized using the Adam optimizer \cite{kingma2015adam} with a learning rate of $5 \times 10^{-4}$. The training process was performed for 500 epochs using Mean Squared Error (MSE) as the loss function:
\begin{equation}
MSE=\frac{1}{N}\sum_{i=1}^{N}(y_i-\hat{y}_i)^2
\end{equation}
where $y_i$ represents the reference solution generated by the numerical simulation, and $\hat{y}_i$ represents the prediction produced by the neural network.


The dataset was divided into training and testing subsets using an 80/20 split. To ensure reproducibility, a fixed random seed was used throughout all experiments, unless otherwise indicated.


\subsection{Dataset Generation}
The datasets used in this work were generated through numerical simulations of different dynamical systems. Each system was simulated independently, producing time-dependent trajectories that were later used as training and evaluation data for the neural network model.

For each system, the input variable was the time coordinate $(t)$, while the target variables corresponded to the physical state of the system, either position or angular displacement.

The generated datasets consisted of trajectories sampled uniformly over a fixed simulation interval. Although the simulation duration and number of samples varied between experiments, each trajectory generally contained between 1,000 and 3,000 samples, allowing the neural network to learn the temporal evolution of the physical process.

\subsubsection{Harmonic Oscillator}
The harmonic oscillator was used as the simplest dynamical system investigated in this study. Its motion is described by the differential equation:
\begin{equation}
\ddot{x}+\omega^2 x=0
\end{equation}

where $x$ represents the displacement and $\omega$ represents the angular frequency.

The analytical solution of this equation is:
\begin{equation}
x(t)=A\cos(\omega t+\phi) 
\end{equation}

where $A$ is the amplitude and $\phi$ is the initial phase.
The generated dataset consists of pairs $(t,x(t))$, where the neural network was trained to approximate the relationship between time and displacement.

\subsubsection{Damped Harmonic Oscillator}
To introduce dissipative effects, a damping term was added to the harmonic oscillator equation:
\begin{equation}
\ddot{x}+\beta\dot{x}+\omega^2 x=0
\end{equation}

where $\beta$ represents the damping coefficient.
Unlike the ideal harmonic oscillator, the amplitude decreases over time due to energy dissipation. The generated trajectories were obtained by numerically solving the differential equation and storing the resulting displacement evolution.

The dataset consisted of time-position pairs, allowing the neural network to learn the effect of damping on the behaviour of the physical system.

\subsubsection{Forced Harmonic Oscillator}
The forced harmonic oscillator extends the damped harmonic oscillator by introducing a periodic external force. This excitation continuously supplies energy to the system, resulting in more complex dynamical behaviour that depends on both the natural frequency of the oscillator and the frequency of the external force.

The motion of the system is governed by:
\begin{equation}
    \ddot{x}+\beta \dot{x}+\omega_0^2x=F_0\sin(\Omega t)
\end{equation}

where $F_0$ represents the forcing amplitude and $\Omega$ represents the forcing angular frequency. The parameters $\beta$ and $\omega_0$ correspond to the damping coefficient and the natural angular frequency, respectively.

Unlike the previous systems, the response of the forced oscillator is determined not only by its intrinsic properties but also by the characteristics of the external excitation. Depending on the relationship between the natural and forcing frequencies, the system may exhibit resonance \cite{landau1986mechanics}, where the oscillation amplitude increases significantly.

As in previous experiments, the trajectories were generated by numerically solving the governing differential equation over a fixed simulation interval. The resulting dataset consisted of pairs $(t,x(t))$, where the input corresponds to the time coordinate and the output represents the corresponding displacement.

Several experiments were performed by varying the forcing amplitude and forcing frequency while keeping the remaining physical parameters constant. This allowed the influence of external excitations on the neural network's ability to approximate the physical system to be evaluated.

\subsubsection{Duffing Oscillator}
The Duffing oscillator was introduced to evaluate the performance of the neural network on a nonlinear dynamical system. Unlike harmonic oscillators, where the restoring force is proportional to the displacement, the Duffing oscillator includes a nonlinear restoring term that can produce complex behaviors depending on the system parameters \cite{strogatz2015nonlinear}.

The governing equation of the Duffing oscillator is
\begin{equation}
    \ddot{x}+\delta\dot{x}+\alpha x+\beta x^3=\gamma\cos(\Omega t),
    \label{eq:duffing_oscillator}
\end{equation}

where $\delta$ represents the damping coefficient, $\alpha$ and $\beta$ control the linear and nonlinear stiffness terms, $\gamma$ is the forcing amplitude, and $\Omega$ is the forcing angular frequency.
The nonlinear term $ (\beta x^3) $ introduces a dependence between the restoring force and the amplitude of the oscillation. Depending on the chosen parameters, the system can exhibit different dynamic regimes, including periodic oscillations, amplitude-dependent frequencies, and more complex responses.

The trajectories were generated by numerically integrating the Duffing differential equation for a predefined set of initial conditions and physical parameters. Each simulation produced a time-dependent displacement trajectory $(x(t))$, which was used to construct the dataset. The neural network received the time coordinate $t$ as input and learned to approximate the corresponding system state.

To analyze the influence of the nonlinear dynamics, experiments were performed by varying selected Duffing parameters while maintaining the remaining parameters fixed. In particular, the damping coefficient $(\delta)$ was modified to study how changes in energy dissipation affected the approximation accuracy of the neural network.

The Duffing oscillator represents a transition between classical oscillatory systems and more complex nonlinear dynamics, providing a challenging benchmark for evaluating whether neural networks can learn physical behaviors beyond simple analytical solutions.

\subsubsection{Simple Pendulum}
The simple pendulum was used as a nonlinear mechanical system to evaluate the ability of the neural network to approximate the dynamics of a physical system with angular motion. Unlike harmonic oscillators, the restoring force of a pendulum depends on the sine of the angular displacement, introducing nonlinear behavior for large oscillations.

The equation of motion of the pendulum is given by
\begin{equation}
    \ddot{\theta}+\frac{g}{L}\sin(\theta)=0,
    \label{eq:simple_pendulum}
\end{equation}

where $\theta$ represents the angular displacement, $g$ is the gravitational acceleration, and $L$ is the pendulum length.
For small angular displacements, the approximation $(\sin(\theta)\approx\theta)$ transforms the system into a harmonic oscillator \cite{goldstein2002classical}. However, in this work the complete nonlinear equation was used, allowing the neural network to learn the full pendulum dynamics without simplifying assumptions.

The trajectories were generated by numerically solving the equation of motion for different initial conditions and physical parameters. Each simulation produced the evolution of the angular position $(\theta(t))$, which was used as the target variable of the neural network. Therefore, the generated dataset consisted of pairs $(t,\theta(t))$, where the input was the time coordinate and the output was the corresponding angular displacement.

Several experiments were performed by varying the gravitational acceleration, pendulum length, and initial angular conditions. These parameter sweeps allowed the analysis of how changes in the physical configuration affected the difficulty of the learning problem and the approximation accuracy of the neural network.

The simple pendulum provides an intermediate benchmark between analytically solvable oscillatory systems and more complex nonlinear systems, as it introduces nonlinear dynamics while maintaining a single degree of freedom.
\subsubsection{Double Pendulum}
The double pendulum was selected as the most complex dynamical system in this study. Unlike the previous systems, which contain a single state variable, the double pendulum is a coupled system with two interacting degrees of freedom. The motion of each pendulum depends not only on its own state but also on the position and velocity of the other pendulum.

The system consists of two pendulum masses connected sequentially, with angular coordinates $\theta_1$ and $\theta_2$. Its dynamics are governed by a coupled set of nonlinear differential equations derived from Lagrangian mechanics \cite{goldstein2002classical}. These equations depend on the masses $m_1$ and $m_2$, the lengths $L_1$ and $L_2$, the gravitational acceleration $g$, and the initial conditions of the system.

The state of the system can be represented as
\begin{equation}
    y(t)=[\theta_1(t),\theta_2(t)],
\end{equation}

where $\theta_1$ and $\theta_2$ describe the angular positions of the first and second pendulum links, respectively.

The trajectories were generated through numerical simulation of the double pendulum equations of motion. For each simulation, the resulting angular trajectories were sampled over a fixed time interval, producing datasets composed of time-state pairs:
\begin{equation}
    (t,[\theta_1(t),\theta_2(t)]).
\end{equation}

In contrast with previous experiments, the neural network was required to approximate a multidimensional output, predicting both angular coordinates simultaneously.

Due to the nonlinear coupling between the two pendulums, small variations in the initial conditions can lead to significantly different trajectories over time. This sensitivity to initial conditions is a characteristic feature of chaotic systems \cite{strogatz2015nonlinear} and makes the double pendulum a challenging benchmark for data-driven modeling approaches.

Experiments were performed by modifying physical parameters such as pendulum lengths and initial angular conditions. These parameter variations were used to evaluate how changes in the underlying dynamics affected the neural network's approximation capabilities.

The double pendulum represents the highest level of complexity considered in this work, extending the evaluation from simple periodic systems to nonlinear and chaotic dynamics.
